(1) 质量线密度 \(\mu>0\)、恒张力 \(T>0\) 的柔弦，左端固定在 \(x=0\)，右端夹持于一个光滑移动的刚性边界 \(x=L(t)\)。小横位移 \(u(x,t)\) 满足波动方程 \(\mu u_{tt}=Tu_{xx}\) 及 \(u(0,t)=u(L(t),t)=0\)。不预设慢变，严格求时变区间的弦能量
\(E(t)=\frac12\int_0^{L(t)}[\mu u_t^2+Tu_x^2]\,dx\) 的时间导数，化成移动端点局部斜率与 \(\dot L\) 的表达式。
(2) 现令 \(|\dot L|/c\ll1\)，其中 \(c=\sqrt{T/\mu}\)，并且模态变化率小于相邻频率差，忽略模间跃迁。求第 \(n\) 瞬时驻波频率，证明单模绝热作用量 \(J_n=E_n/\omega_n\) 不变，推导单模能量随 \(L\) 的标度及腹部横向位移振幅的标度。
(3) 对弦位移近似 \(u=A\sin(n\pi x/L)\cos\omega_nt\)，独立用第一问的能量边界式求慢变平均夹持广义力 \(\langle F_{\rm end}\rangle\) 与 \(\dot E\)。与第二问预测 \(-dE/dL\) 比较，解释正向伸长为何使单模能量下降。
(4) 换一种独立机制：考虑速度 \(v_w=\dot L\) 缓慢远离入射波的反射端点，推导实验室中单次一维波反射的频率比 \(\omega_r/\omega_i=(c-v_w)/(c+v_w)\)。结合两端往返撞击率 \(c/(2L)\) 恢复 \(\dot E/E\simeq-\dot L/L\)，说明非绝热快速移动时此估计为什么失效。
